% Front matter of the edition (numbered separately; the text keeps the lithograph's page numbers)
\pagenumbering{alph}
\thispagestyle{frontframe}
\begin{center}
\vspace*{2mm}
{\large كتاب\par}
\vspace{4mm}
\scalebox{0.95}{\wafiyahtitle}\par
\vspace{5mm}
{\large في الصرف\par}
\vspace{7mm}
{\normalsize تأليف\par}
\vspace{2mm}
{\Large\bfseries الشيخ محمد ميران الباقوي القاسمي\par}
\vspace{1mm}
{\small مدرّس دار العلوم سبيل الرشاد، بنغلور\par}
\vspace{7mm}
\ornament\par
\vspace{7mm}
{\normalsize دراسة وتحقيق\par}
\vspace{2mm}
{\Large\bfseries\textcolor{gray}{[اسم المحقق]}\par}
\vfill
{\small عن الطبعة الحجرية: سبيل الفلاح، حيدرآباد، ١٣٩٦هـ / ١٩٧٦م\par}
\end{center}
\clearpage

% Key to the editorial signs used in the text and the apparatus
\thispagestyle{frontframe}
\gdef\currentbab{}
\vspace*{\fill}
\begin{center}{\large\bfseries رموز التحقيق ومصطلحاته\par}\vspace{2mm}\ornament\end{center}
\vspace{5mm}
\begin{center}
\renewcommand{\arraystretch}{1.8}
\begin{tabular}{@{}c@{\hspace{5mm}}r@{}}
\textsuperscript{\latinfont†} & قراءة غير واضحة في الأصل، أُثبتت كما قُرئت\\
\textcolor{gray}{[\,…\,]} & بياض أو كلمة لم تُقرأ في الأصل\\
\x & لا صيغة لهذا الموضع، كما في الأصل\\
\raisebox{.3ex}{\textcolor{gray}{\latinfont\small″}} & علامة التكرار: مثل ما فوقه\\
{\footnotesize الخ} & إلى آخره\\
\colnum{5} & رقم العمود في الأصل\\
\end{tabular}
\end{center}
\vspace{6mm}
\begin{center}
\begin{minipage}{86mm}\small\setlength{\parskip}{4pt}
\textbf{الحواشي:} الحاشية الأولى لما في الأصل من قراءات غير واضحة وتصحيحات، وتُذكر فيها اللفظة ثم علامة \mbox{(\,]\,)} ثم البيان؛ والحاشية الثانية لتعليقات المحقق.

\textbf{أرقام الصفحات:} تطابق أرقامُ صفحات هذه النشرة أرقامَ صفحات الطبعة الحجرية، كلُّ صفحة منها في صفحة.
\end{minipage}
\end{center}
\vspace*{\fill}
\clearpage
\pagenumbering{arabic}
